\section{Apêndice A -- Premissas e fórmulas da simulação}
\label{ap:modelo}

A simulação reproduz, por fórmula, os resultados do mês 12 da Atividade IX (receita de R\$~4.421,50; margem de contribuição de R\$~2.232,15; resultado do cenário A de R\$~756,32 e do cenário B de R\$~$-$4.743,68). Seja $N$ o número de usuários ativos mensais.

\begin{equation}
\text{Receita}(N) = N\,[\,a\,t\,c\,k + s\,p + \rho\,m\,] + C
\end{equation}

em que $a$ é a proporção de usuários com atendimento intermediado, $t$ o ticket médio, $c$ a parcela paga em dinheiro, $k$ a comissão, $s$ a proporção de assinantes, $p$ o preço da assinatura, $\rho$ o número de clínicas pagantes por usuário, $m$ a mensalidade do Plano Parceiro e $C$ o convênio mensal. As deduções são:

\begin{equation}
\text{Taxas}(N) = N\,[\,s\,p\,\lambda + a\,t\,c\,\phi + \rho\,m\,\psi\,]
\end{equation}
\begin{equation}
\text{Contribuição}(N) = \text{Receita}(N) - \text{Taxas}(N) - \tau\,\text{Receita}(N) - N\,(i + a\,f)
\end{equation}

em que $\lambda$ é a taxa da loja de aplicativos, $\phi$ a taxa média de pagamento, $\psi$ a taxa de cartão sobre o Plano Parceiro, $\tau$ a alíquota de impostos, $i$ a infraestrutura e o suporte por usuário e $f$ o fundo de liquidação por atendimento. O resultado mensal é a contribuição menos os custos fixos do cenário A (R\$~1.475,83) ou B (R\$~6.975,83).

\begin{table}[htb]
\centering
\caption{Premissas do cenário-base}
\label{tab:premissas}
\fontsizetable
\begin{tabularx}{\textwidth}{Y L{3.2cm}}
\toprule
\textbf{Parâmetro} & \textbf{Valor}\\
\midrule
Atendimentos intermediados ($a$) & 10\% dos usuários\\
Ticket médio ($t$) e parcela em dinheiro ($c$) & R\$~110,00; 70\%\\
Comissão ($k$) & 10\%\\
Assinantes ($s$) e preço ($p$) & 5\%; R\$~12,90\\
Clínicas pagantes ($\rho$) e mensalidade ($m$) & 1 por 150 usuários; R\$~79,90\\
Convênio ($C$) & R\$~1.500,00\\
Taxas: loja ($\lambda$), pagamento ($\phi$), cartão ($\psi$) & 15\%; 2,59\%; 4,98\%\\
Impostos ($\tau$) & 6\%\\
Infraestrutura e suporte ($i$) & R\$~0,30\\
Fundo de liquidação ($f$) & R\$~6,60 por atendimento\\
Custos fixos (A; B) & R\$~1.475,83; R\$~6.975,83\\
\bottomrule
\end{tabularx}
\fontetexto{Atividade IX \cite{grupo2026a9}}
\end{table}

De forma simplificada, o ponto de equilíbrio pode ser lido como $N_{eq} = (C_f - R_{fixa})/m_u$, em que $C_f$ são os custos fixos mensais, $R_{fixa}$ a receita líquida que independe do número de usuários (por exemplo, o convênio, após impostos) e $m_u$ a contribuição por usuário. Essas expressões são instrumentos de análise a atualizar com dados observados no piloto, e não projeções ou números já medidos.

Os pontos de equilíbrio foram obtidos por busca numérica da menor quantidade de usuários para a qual o resultado é não negativo, arredondada para cima. Para o cenário-base, obtêm-se 121 usuários (A, com convênio), 2.693 (A, sem convênio), 10.155 (B, com convênio) e 12.728 (B, sem convênio), valores que coincidem com os da Atividade IX, exceto por diferenças de arredondamento. As alavancas alteram, de forma isolada ou conjunta, $f$, $s$, $\rho$ e $p$.
